\documentclass{article}
\usepackage{graphicx} % Required for inserting images
\usepackage[english]{babel}
\usepackage{amsmath}
\usepackage{amsthm}
\usepackage{amssymb}
\usepackage[margin=.75in]{geometry}
\usepackage{hyperref}
\usepackage{amsfonts}
\usepackage{mathrsfs}
\usepackage[version=4]{mhchem}
\usepackage{bigints}
\usepackage{caption}
\usepackage{xpatch}
\usepackage{xr}
\usepackage{enumitem}

\title{Math 330 Abstract Algebra Notes}
\author{Kyle Wang}
\date{September 2026}

\newcommand{\R}{\mathbb{R}}
\newcommand{\Z}{\mathbb{Z}}
\newcommand{\N}{\mathbb{N}}

\newtheorem{theorem}{Theorem}[section]
\newtheorem{lemma}[theorem]{Lemma}
\newtheorem{corollary}[theorem]{Corollary}
\theoremstyle{definition}
\newtheorem{definition}[theorem]{Definition}
\newtheorem{example}[theorem]{Example}

\begin{document}
\maketitle
\tableofcontents
\section{Group Theory}
\subsection{Modular Arithmetic}
\begin{definition}
    Take $n>0$ and $n\in \Z $. We say that $a,b\in \Z$ are \textit{congruent modulo n} ($a\equiv b \mod{n}$) if $n$ divides $a-b$ ($n|a-b$) or if the remainder of $n/a$ equals the remainder of $n/b$.
\end{definition}
Here are some properties:
\begin{enumerate}
    \item if $a\equiv b \mod n$ and $c\equiv d \mod n$, then $a+c\equiv b+d \mod n$.
    \begin{proof}
        $a\equiv b \mod n$ implies that $n|a-b$ which implies that $\exists q \in \Z$ such that $a-b=nq$. The same can be said for $a\equiv b \mod n$ for some $q' \in \Z$ such that $c-d=nq'$. Adding these equations together we have that $a-b+c-d = (a+c)-(b+d) = n(q+q')$. Since the integers are closed under addition, $q+q'$ is an integer and therefore $n|(a+c)-(b+d)$ and as a result, $a+c\equiv b+d \mod n$.
    \end{proof}
    \item if $a\equiv b \mod n$ and $c\equiv d \mod n$, then $ac\equiv bd \mod n$
    \item Congruence $\!\!\!\!\mod n$ is an equivalence relation
\end{enumerate}
\subsection{Introduction to Groups}
\begin{definition}[Group]
    A group is a set of elements $G$ with a binary operation $*$ that assigns to every pair $(g,h)$ of elements of $G$ another element in $G$ denoted $g*h$. Groups follow the group axioms:
    \begin{enumerate}
        \item Associativity
        \begin{equation*}
            (g_1*g_2)*g_3 = g_1*(g_2*g_3)
        \end{equation*}
        \item There exists an element $e\in G$ that is the identity such that 
        \begin{equation*}
            e*g_1=g_1*e=g_1
        \end{equation*}
        \item Every element has an inverse element denoted $g^{-1}$ such that 
        \begin{equation*}
            g*g^{-1} = g^{-1}*g = e
        \end{equation*}
        \item The set of elements is closed under the group operation (although this is implied in the definition).
    \end{enumerate}
\end{definition}
\begin{example}
    Let us look at some examples
    \begin{enumerate}
        \item The set $(\Z,+)$. Addition is associative, the identity is 0, the inverse of each element $g$ is $-g$. 
        \item The set $(\R_{>0}, \cdot)$ where $\cdot$ denotes regular multiplication. Regular multiplication is associative, the identity is $1$, the inverse element is $\frac{1}{g}$
        \item The set $(\{-2,-1,0,1,2\},+)$. This fails due since it is not closed (take $1+2=3 \notin $ the set
        \item The set $\left(\left\{ \frac{1}{2^n} \right\}, \cdot \right)$. This fails since there is no inverse.
        \item Fix a positive integer $n$. Take the set $\left( \{0,1,2,\ldots,n-1\}, +\mod n\right)$. The identity is 0, addition is associative. The inverse of an element $0\leq a\leq n-1$ is $n-a$ since $(n-a)+a \mod n = n\mod n = 0$. This group is denoted by Fraleigh as $\Z$ but more generally denoted as $\Z/n\Z$.
        \item Fix an $n\in \N$. We denote the group GL$_n(\R)$ as the set of invertible $n\times n$ matrices with real entries with the operation of matrix multiplication. Matrix multiplication is associative, the identity is the $n\times n$ identity matrix $I_n$. The inverse of every element $A$ is the inverse matrix of $A$. This group is also known as the \textit{General Linear Group}.
    \end{enumerate}
\end{example}

\end{document}
